\documentclass[10pt,letterpaper]{article}

\usepackage[top=0.8in,bottom=0.8in,left=0.9in,right=0.9in]{geometry}
\usepackage[default,scale=1.0]{lato}
\usepackage[T1]{fontenc}
\usepackage{enumitem}
\usepackage{xcolor}
\usepackage[colorlinks=true,urlcolor=linkteal,linkcolor=linkteal]{hyperref}

\definecolor{linkteal}{HTML}{00707A}
\definecolor{graytext}{HTML}{666666}

\setlength{\parindent}{0pt}
\pagestyle{empty}

% ---- 見出し：太字 + 全幅の罫線 ----
\newcommand{\cvsection}[1]{%
  \vspace{1.05em}%
  {\fontsize{11.5}{14}\selectfont\bfseries #1}\par
  \vspace{3pt}%
  \hrule height 0.5pt
  \vspace{0.62em}%
}

% ---- 本文：字下げしたラベルなしリスト ----
\newenvironment{cvlist}{%
  \begin{list}{}{%
    \setlength{\leftmargin}{1.7em}%
    \setlength{\labelwidth}{0pt}%
    \setlength{\labelsep}{0pt}%
    \setlength{\itemsep}{0.50em}%
    \setlength{\parsep}{0pt}%
    \setlength{\topsep}{0pt}%
    \setlength{\partopsep}{0pt}%
  }%
}{\end{list}}

% ---- 行間を詰めた版（学歴・一覧向け）----
\newenvironment{cvlisttight}{%
  \begin{list}{}{%
    \setlength{\leftmargin}{1.7em}%
    \setlength{\labelwidth}{0pt}%
    \setlength{\labelsep}{0pt}%
    \setlength{\itemsep}{0.32em}%
    \setlength{\parsep}{0pt}%
    \setlength{\topsep}{0pt}%
    \setlength{\partopsep}{0pt}%
  }%
}{\end{list}}

\newcommand{\entry}{\item[]}

% ---- 要旨：一段下げ・小さめ・両端揃え ----
\newenvironment{abstractblock}{%
  \begingroup
  \fontsize{8.6}{11.3}\selectfont
  \begin{list}{}{%
    \setlength{\leftmargin}{0.9em}%
    \setlength{\rightmargin}{0.4em}%
    \setlength{\topsep}{0.15em}%
    \setlength{\parsep}{0pt}%
    \setlength{\itemsep}{0pt}%
  }\item[]\textbf{Abstract.}\ %
}{\end{list}\endgroup}
\newcommand{\dated}[2]{#1\hfill\textit{#2}}
\newcommand{\subhead}[1]{\item[]{\bfseries #1}}

\begin{document}

\begin{center}
{\LARGE\bfseries Hideyuki Tomiyama}
\end{center}
\vspace{0.6em}

\cvsection{Contact Information}
\begin{cvlisttight}
\entry Email: \href{mailto:hkt5140@psu.edu}{hkt5140@psu.edu}
\entry Website: \href{https://thideyuki.github.io}{\texttt{https://thideyuki.github.io}}
\end{cvlisttight}

\cvsection{Education}
\begin{cvlist}
\entry Ph.D.\ in Economics, The Pennsylvania State University, 2021 -- 2027 (expected)
\entry M.A.\ in Economics, Keio University, 2018 -- 2020
\entry B.A.\ in Economics, Keio University, 2014 -- 2018
\end{cvlist}

\cvsection{Fields of Interest}
\begin{cvlisttight}
\entry Econometrics, Labor Economics, Causal Inference, Applied Econometrics
\end{cvlisttight}

\cvsection{Job Market Paper}
\begin{cvlist}
\entry Causal Model of a Two-Sided Matching Market: Application to Intergenerational
       Spillover Effects of Education
\item[]\vspace{-0.15em}
\begin{abstractblock}
This paper studies a heterogeneous treatment effect model with instrumental variables in an endogenous one-to-one two-sided matching market. Our leading example is the effect of parental education on children's outcomes, such as income, educational attainment, and health. Under assortative matching, it is difficult to justify the standard instrumental variables methods (e.g., two-stage least squares) that use treatments and instruments from both sides. To overcome this problem, we explicitly consider a selection model in which each agent selects both their own treatment status and that of their matched partner. We provide point and partial identification results for various causal parameters, including the average treatment effect. The key identification assumption is matching on observables, which requires that pairs are randomly matched conditional on observed characteristics and rules out assortative matching on unobservables. We also propose inference procedures for these causal parameters. Finally, we apply our framework to study the treatment effect of parents' college attendance on the educational investment they make in their child, using Japanese household-level data.
\end{abstractblock}
\end{cvlist}

\cvsection{Working Papers}
\begin{cvlist}
\entry Identification and Estimation of Panel Regression Models with Unknown Aggregated
       Spatial Units
\item[]\vspace{-0.15em}
\begin{abstractblock}
This paper studies panel regression models in which each unit corresponds to an aggregated spatial unit. Each aggregated unit is a collection of multiple finer units in the raw dataset. The variables of each aggregated unit can be constructed as the sum of finer unit-level variables in the aggregated unit. However, the researcher does not know the correct delineation of aggregated units under which the structural equation holds. Examples include regression models of agglomeration economies and game theoretic models of spatial entry competitions. We show that the derivative of the structural function can be nonparametrically identified from observational panel data on finer units under certain conditions. The key identification condition is the mean independence between the structural error terms and the vector of finer unit-level exogenous variables. The correct delineation can be also identified if the structural function is nonaffine with respect to the regressors. The framework can be applied to other aggregation problems such as industry aggregation. We propose several computationally feasible inference procedures for parametric models via GMM-type objective functions. As an empirical illustration, we apply the methodology to the economics of agglomeration using U.S. data.
\end{abstractblock}
\end{cvlist}

\cvsection{Publications}
\begin{cvlist}
\entry Inference on Incomplete Information Games with Multi-Dimensional Actions,
       with Taisuke Otsu. \textit{Economics Letters}, Volume 215, 2022
\end{cvlist}

\cvsection{Teaching Experience}
\begin{cvlisttight}
\subhead{The Pennsylvania State University (Teaching Assistant)}
\entry \dated{ECON 402, Decision Making and Strategy in Economics}{Spring 2026}
\entry \dated{ECON 454, Economics of Mergers}{Fall 2024 -- Fall 2025}
\entry \dated{ECON 315, Labor Economics}{Spring 2024}
\entry \dated{ECON 451, Monetary Theory and Policy}{Fall 2023}
\entry \dated{ECON 104, Introductory Macroeconomic Analysis and Policy}{Spring 2022}
\entry \dated{ECON 102, Introductory Microeconomic Analysis and Policy}{Fall 2021}
\vspace{0.45em}
\subhead{Keio University (Teaching Assistant)}
\entry \dated{Industrial Organization, Economic Development,
       Intermediate Macroeconomics}{2018 -- 2021}
\end{cvlisttight}

\cvsection{Seminars and Conferences}
\begin{cvlist}
\entry 2026: Summer Workshop on Economic Theory;
       The Pennsylvania State University (scheduled)
\entry 2020: Summer Workshop on Economic Theory;
       Japanese Economic Association Spring Meeting
\end{cvlist}

\cvsection{Fellowships}
\begin{cvlist}
\entry 2021 -- present: Graduate Assistantship,
       The Pennsylvania State University
\entry 2020 -- 2021: Research Fellowship for Young Scientists (DC1),
       Japan Society for the Promotion of Science
\entry 2020 -- 2021: Adjunct Researcher, Faculty of Economics, Keio University
\end{cvlist}

\cvsection{Languages}
\begin{cvlisttight}
\entry English (fluent), Japanese (native)
\end{cvlisttight}

\cvsection{Citizenship}
\begin{cvlisttight}
\entry Japan
\end{cvlisttight}

\cvsection{References}
\begin{cvlisttight}
\item[]
\begin{minipage}[t]{0.31\linewidth}
{\bfseries Professor Marc Henry}\\
Department of Economics,\\
Pennsylvania State University\\
Email: marc.henry@psu.edu\\
Phone: +1(814)865-0010
\end{minipage}\hfill
\begin{minipage}[t]{0.31\linewidth}
{\bfseries Professor Keisuke Hirano}\\
Department of Economics,\\
Pennsylvania State University\\
Email: kuh237@psu.edu\\
Phone: +1(814)867-3312
\end{minipage}\hfill
\begin{minipage}[t]{0.31\linewidth}
{\bfseries Professor Sung Jae Jun}\\
Department of Economics,\\
Pennsylvania State University\\
Email: suj14@psu.edu\\
Phone: +1(814)865-6149
\end{minipage}
\end{cvlisttight}

\end{document}
